\documentclass[10pt,a4paper]{article}
\usepackage{webclases-lesson}
\hypersetup{pdftitle={T0 - Herramientas: resumen en una pagina}, pdfauthor={Web Clases Rocedg}}

% Ajustes propios del resumen: A4 apaisado, sin pie y cajas compactas.
\geometry{landscape,left=7mm,right=7mm,top=6mm,bottom=6mm}
\pagestyle{empty}
\linespread{1.0}
\setlength{\parskip}{0.15em}
\setlength{\tabcolsep}{0.3em}
\renewcommand{\arraystretch}{1.12}
\setlength{\abovedisplayskip}{2pt}
\setlength{\belowdisplayskip}{2pt}
\setlength{\abovedisplayshortskip}{1pt}
\setlength{\belowdisplayshortskip}{2pt}
\setlist[itemize]{leftmargin=0.9em,topsep=0.05em,itemsep=0em,parsep=0em,label=\textcolor{Muted}{\textbullet}}

% Un color por grupo: calculo (azul), vectores (violeta), unidades (verde), medida (ambar).
\tcbset{sumbox/.style n args={3}{enhanced,colback=#1,colframe=#2!32,
  borderline west={2.2pt}{0pt}{#2},boxrule=0.35pt,arc=2mm,
  left=2.0mm,right=1.6mm,top=0.9mm,bottom=1.0mm,before skip=0pt,after skip=1.5mm,
  fontupper=\small,
  before upper={\raggedright\textbf{\textcolor{#2!80!black}{#3}}\par\vspace{0.15em}}}}
\newtcolorbox{calcbox}[1]{sumbox={BlueSoft}{BrandBlue}{#1}}
\newtcolorbox{vecbox}[1]{sumbox={ToolSoft}{ToolLine}{#1}}
\newtcolorbox{unitbox}[1]{sumbox={GreenSoft}{GreenLine}{#1}}
\newtcolorbox{measbox}[1]{sumbox={AmberSoft}{AmberLine}{#1}}
\newcommand{\grouptag}[3]{\tcbox[on line,colback=#1,colframe=#2!40,boxrule=0.35pt,arc=1.5mm,
  left=1.2mm,right=1.2mm,top=0.3mm,bottom=0.3mm]{\footnotesize\bfseries\textcolor{#2!80!black}{#3}}}
\newcommand{\lbl}[1]{\textbf{\textcolor{Ink}{#1}}\enspace}
\newcommand{\minus}{\mathbin{\color{RedLine}\boldsymbol{-}}}

\begin{document}
\shorthandoff{<>}
\color{Ink}

% Cabecera
{\Large\bfseries\textcolor{Ink}{T0 \textperiodcentered\ Herramientas: resumen en una p\'agina}}%
\hfill
\grouptag{BlueSoft}{BrandBlue}{C\'alculo}\,%
\grouptag{ToolSoft}{ToolLine}{Vectores}\,%
\grouptag{GreenSoft}{GreenLine}{Unidades SI}\,%
\grouptag{AmberSoft}{AmberLine}{Medida}\quad
{\small\bfseries\textcolor{BrandBlueDark}{Lecci\'on T0 \textperiodcentered\ v1.0}}\enspace
{\small\textcolor{Muted}{1.\textordmasculine\ Bachillerato}}\par
\vspace{0.2em}{\color{BrandBlue}\hrule height 0.8pt}\vspace{0.45em}

\noindent
% ===================== Columna 1: calculo =====================
\begin{minipage}[t]{0.322\linewidth}
\vspace{0pt}
\begin{calcbox}{1 \textperiodcentered\ Derivadas}
\renewcommand{\arraystretch}{1.3}%
\hspace*{0.5em}%
\begin{tabular}[t]{@{}>{$}r<{$}@{$\ \to\ $}>{$}l<{$}@{}}
c & 0\\
x^n & nx^{n-1}\\
\sqrt x & \dfrac{1}{2\sqrt x}\\
\dfrac1x & -\dfrac{1}{x^2}\\
\end{tabular}\hfill
\begin{tabular}[t]{@{}>{$}r<{$}@{$\ \to\ $}>{$}l<{$}@{}}
\sin x & \cos x\\
\cos x & -\sin x\\
e^x & e^x\\
\end{tabular}\hspace*{1.5em}%
\renewcommand{\arraystretch}{1.12}%
\par\vspace{0.25em}
\lbl{Reglas}
$(c\,f)'=c\,f'$\quad $(f+g)'=f'+g'$\quad $(fg)'=f'g+fg'$\par
$\left(\dfrac fg\right)'=\dfrac{f'g-fg'}{g^2}$\qquad
\lbl{Cadena} $\big[f(g(x))\big]'=f'(g(x))\cdot g'(x)$\par
\lbl{Ejemplo} $\dfrac{d}{dt}\sin(\omega t)=\omega\cos(\omega t)$
\end{calcbox}

\begin{calcbox}{2 \textperiodcentered\ Integrales}
\renewcommand{\arraystretch}{1.5}%
\begin{tabular}[t]{@{}>{$\displaystyle}l<{$}@{}}
\int x^n\,dx=\frac{x^{n+1}}{n+1}+C\\[-3pt]
\qquad\textcolor{RedLine!80!black}{(n\ne-1)}\\
\int\frac1x\,dx=\ln|x|+C\\
\end{tabular}\hfill
\begin{tabular}[t]{@{}>{$\displaystyle}l<{$}@{}}
\int\sin x\,dx=-\cos x+C\\
\int\cos x\,dx=\sin x+C\\
\int e^x\,dx=e^x+C\\
\end{tabular}%
\renewcommand{\arraystretch}{1.12}%
\par\vspace{0.25em}
\lbl{Definida} $\displaystyle\int_a^b f\,dx=F(b)-F(a)$, con $F'=f$.\par
\lbl{La constante $C$} la fija una condici\'on inicial:
$v=3+2t\Rightarrow x=3t+t^2+C$; si $x(0)=2\units{m}$, $C=2\units{m}$.
\end{calcbox}

\begin{calcbox}{3 \textperiodcentered\ Derivada e integral en F\'isica}
\centering
\begin{tikzpicture}[font=\small,
  q/.style={draw=BrandBlue!45,fill=white,rounded corners=1.5mm,minimum width=7mm,minimum height=5mm}]
\node[q] (x) at (0,0) {$x$};
\node[q] (v) at (2.6,0) {$v$};
\node[q] (a) at (5.2,0) {$a$};
\draw[-{Latex[length=1.8mm]},BrandBlue,thick] ([yshift=1mm]x.east)--([yshift=1mm]v.west) node[midway,above] {$\frac{d}{dt}$};
\draw[-{Latex[length=1.8mm]},BrandBlue,thick] ([yshift=1mm]v.east)--([yshift=1mm]a.west) node[midway,above] {$\frac{d}{dt}$};
\draw[-{Latex[length=1.8mm]},GreenLine!80!black,thick] ([yshift=-1mm]v.west)--([yshift=-1mm]x.east) node[midway,below] {$\int dt$};
\draw[-{Latex[length=1.8mm]},GreenLine!80!black,thick] ([yshift=-1mm]a.west)--([yshift=-1mm]v.east) node[midway,below] {$\int dt$};
\end{tikzpicture}\par
\raggedright
$v=\dfrac{dx}{dt}$,\ $a=\dfrac{dv}{dt}$\qquad
$\Delta x=\displaystyle\int_{t_1}^{t_2}v\,dt$,\ $\Delta v=\displaystyle\int_{t_1}^{t_2}a\,dt$\par
\vspace{0.25em}
\begin{tikzpicture}[x=1.05cm,y=1.05cm,font=\footnotesize]
\draw[axis] (0,0)--(2.75,0) node[below left=-1pt] {$t$};
\draw[axis] (0,0)--(0,2.05) node[left] {$x$};
\draw[BrandBlue,thick,domain=0.05:2.6,samples=40] plot(\x,{0.22*\x*\x+0.25});
\draw[GreenLine!80!black,thick] (0.75,0.3875)--(2.55,1.6475);
\fill[BrandBlue] (1.5,0.745) circle (1.5pt);
\draw[guide] (1.5,0.745)--(2.2,0.745)--(2.2,1.235);
\node[below,inner sep=1pt] at (1.85,0.745) {$\Delta t$};
\node[right,inner sep=1pt] at (2.2,0.99) {$\Delta x$};
\node[align=center,GreenLine!70!black] at (1.15,1.75) {pendiente de\\la tangente $=v$};
\end{tikzpicture}\hfill
\begin{tikzpicture}[x=1.05cm,y=1.05cm,font=\footnotesize]
\fill[BrandBlue!25] (0,0)--(0,1.15)--(1.6,0)--cycle;
\fill[RedLine!25] (1.6,0)--(2.5,-0.646875)--(2.5,0)--cycle;
\draw[axis] (0,-0.85)--(0,1.4) node[left] {$v$};
\draw[axis] (0,0)--(2.8,0) node[above left=-1pt] {$t$};
\draw[BrandBlue,thick] (0,1.15)--(2.5,-0.646875);
\node[BrandBlueDark] at (0.45,0.35) {$+$};
\node[RedLine!80!black] at (2.2,-0.22) {$-$};
\node[align=center,anchor=west] at (0.85,1.05) {\'area $=\Delta x$};
\node[align=left,anchor=north west,RedLine!80!black] at (0.05,-0.18) {bajo el eje:\\negativa};
\end{tikzpicture}\par
\vspace{0.15em}
\lbl{M\'aximo o m\'inimo} $f'=0$ \textbf{y} $f'$ cambia de signo:
$+\to-$ m\'aximo; $-\to+$ m\'inimo. ($x^3$: $f'(0)=0$ sin extremo.)
\end{calcbox}
\end{minipage}\hfill
% ===================== Columna 2: vectores =====================
\begin{minipage}[t]{0.322\linewidth}
\vspace{0pt}
\begin{vecbox}{4 \textperiodcentered\ Formas de escribir un vector}
\begin{minipage}[c]{0.56\linewidth}
\raggedright
\lbl{Par} $\vec A=(A_x;\ A_y)$\par
\lbl{Columna} $\vec A=\begin{pmatrix}A_x\\A_y\end{pmatrix}$\par
\lbl{Unitarios} $\vec A=A_x\ii+A_y\jj$\par
\lbl{Polar} $\vec A=(r\cos\theta;\ r\sin\theta)$
\end{minipage}\hfill
\begin{minipage}[c]{0.42\linewidth}
\centering
\begin{tikzpicture}[x=1cm,y=1cm,font=\footnotesize]
\draw[axis] (-.1,0)--(2.45,0) node[below left=-1pt] {$x$};
\draw[axis] (0,-.1)--(0,1.65) node[left] {$y$};
\draw[guide] (2,0)--(2,1.5)--(0,1.5);
\draw[vector] (0,0)--(2,1.5) node[midway,sloped,above=-1pt] {$r=|\vec A|$};
\draw[-{Latex[length=1.3mm]}] (.6,0) arc[start angle=0,end angle=36.87,radius=.6];
\node at (.8,.22) {$\theta$};
\node[below] at (1,0) {$A_x$};
\node[right] at (2,.75) {$A_y$};
\end{tikzpicture}
\end{minipage}\par\vspace{0.2em}
$A_x=|\vec A|\cos\theta$,\ $A_y=|\vec A|\sin\theta$;\enspace $\theta$ desde $+x$, antihorario.\par
\lbl{M\'odulo} $|\vec A|=\sqrt{A_x^2+A_y^2\,(+A_z^2)}$\qquad
$\tan\theta=\dfrac{A_y}{A_x}$\par
\textcolor{AmberLine!70!black}{\textbf{Ojo al cuadrante:}} mira los signos.
$(-4;\ 3)$: la calculadora da $-36{,}9^\circ$, pero $\theta=143{,}1^\circ$ (2.\textordmasculine\ cuadrante).
\end{vecbox}

\begin{vecbox}{5 \textperiodcentered\ Vector unitario}
$\displaystyle\hat u_A=\frac{\vec A}{|\vec A|}
=\left(\frac{A_x}{|\vec A|};\ \frac{A_y}{|\vec A|}\right)$\quad ($\vec A\ne\vec 0$)\par
\lbl{Comprueba} $|\hat u_A|=1$.\enspace
$(6;\ -8)/10=(0{,}6;\ -0{,}8)$ y $0{,}36+0{,}64=1$.
\end{vecbox}

\begin{vecbox}{6 \textperiodcentered\ Producto escalar}
$\vec A\cdot\vec B=A_xB_x+A_yB_y+A_zB_z=|\vec A||\vec B|\cos\theta$\par
Da un \textbf{n\'umero}.\enspace $\vec A\cdot\vec B=0\iff$ perpendiculares.\enspace
\'Angulo: $\cos\theta=\dfrac{\vec A\cdot\vec B}{|\vec A||\vec B|}$\par
\lbl{Ejemplo} $(4;\ 3)\cdot(-1;\ 2)=-4+6=2$.\enspace Trabajo: $W=\vec F\cdot\Delta\vec r$.
\end{vecbox}

\begin{vecbox}{7 \textperiodcentered\ Producto vectorial}
$\vec A\times\vec B=
\begin{vmatrix}\ii&\jj&\hat k\\A_x&A_y&A_z\\B_x&B_y&B_z\end{vmatrix}$\par\vspace{0.1em}
\lbl{Sarrus} copia las dos primeras columnas a la derecha.\par\vspace{0.15em}
\begin{tikzpicture}[x=0.78cm,y=0.5cm,font=\small]
\foreach \c in {0,1,2} {
  \draw[GreenLine!80!black,thick] ($(\c,2)+(.17,-.22)$)--($(\c+2,0)+(-.17,.22)$);
}
\foreach \c in {2,3,4} {
  \draw[RedLine,thick,dashed] ($(\c,2)+(-.17,-.22)$)--($(\c-2,0)+(.17,.22)$);
}
\draw[thick] (-.4,-.35)--(-.4,2.35) (2.4,-.35)--(2.4,2.35);
\tikzset{entry/.style={fill=ToolSoft,inner sep=1pt}}
\node[entry] at (0,2) {$\ii$}; \node[entry] at (1,2) {$\jj$}; \node[entry] at (2,2) {$\hat k$};
\node[entry,text=Muted] at (3,2) {$\ii$}; \node[entry,text=Muted] at (4,2) {$\jj$};
\node[entry] at (0,1) {$A_x$}; \node[entry] at (1,1) {$A_y$}; \node[entry] at (2,1) {$A_z$};
\node[entry,text=Muted] at (3,1) {$A_x$}; \node[entry,text=Muted] at (4,1) {$A_y$};
\node[entry] at (0,0) {$B_x$}; \node[entry] at (1,0) {$B_y$}; \node[entry] at (2,0) {$B_z$};
\node[entry,text=Muted] at (3,0) {$B_x$}; \node[entry,text=Muted] at (4,0) {$B_y$};
\node[right,align=left,font=\footnotesize] at (4.5,1)
  {\textcolor{GreenLine!70!black}{continua ($\searrow$): suman}\\[1pt]
   \textcolor{RedLine}{discontinua ($\swarrow$): restan}\\[1pt]
   \textcolor{Muted}{gris: columnas copiadas}};
\end{tikzpicture}\par\vspace{0.15em}
$\vec A\times\vec B=A_yB_z\,\ii+A_zB_x\,\jj+A_xB_y\,\hat k$\par
\hspace*{3.2em}$-\left(A_yB_x\,\hat k+A_zB_y\,\ii+A_xB_z\,\jj\right)$\par\vspace{0.1em}
\lbl{Componente a componente}\par
\hspace*{0.6em}$=(A_yB_z-A_zB_y)\,\ii+(A_zB_x-A_xB_z)\,\jj+(A_xB_y-A_yB_x)\,\hat k$\par\vspace{0.2em}
\begin{minipage}[c]{0.55\linewidth}
\raggedright
$|\vec A\times\vec B|=|\vec A||\vec B|\sin\theta$ $=$ \'area del paralelogramo.\par
Perpendicular a $\vec A$ y a $\vec B$.\par
$\vec B\times\vec A=-\vec A\times\vec B$\par
$\ii\times\jj=\hat k$;\ $+\hat k$ sale del papel ($\odot$).
\end{minipage}\hfill
\begin{minipage}[c]{0.43\linewidth}
\centering
\begin{tikzpicture}[x=1cm,y=1cm,font=\footnotesize]
\fill[ToolLine!14] (0,0)--(1.7,0)--(2.45,0.6)--(0.75,0.6)--cycle;
\draw[guide] (1.7,0)--(2.45,0.6)--(0.75,0.6);
\draw[vector] (0,0)--(1.7,0) node[below] {$\vec A$};
\draw[vectorgreen] (0,0)--(0.75,0.6) node[left=1pt] {$\vec B$};
\draw[vectorred] (0,0)--(0,1.55) node[left] {$\vec A\times\vec B$};
\draw[-{Latex[length=1.3mm]},AmberLine!85!black,thick] (0.55,0) arc[start angle=0,end angle=38.66,radius=0.55];
\node[AmberLine!70!black,align=left,anchor=west] at (0.2,1.3) {dedos: $\vec A\to\vec B$\\pulgar: $\vec A\times\vec B$};
\end{tikzpicture}
\end{minipage}
\end{vecbox}
\end{minipage}\hfill
% ===================== Columna 3: unidades y medida =====================
\begin{minipage}[t]{0.322\linewidth}
\vspace{0pt}
\begin{unitbox}{8 \textperiodcentered\ Unidades SI}
\lbl{Magnitudes base}\par
\begin{tabular}{@{}l@{\quad}r@{\qquad}l@{\quad}r@{}}
longitud & m & corriente & A\\
masa & kg & temperatura & K\\
tiempo & s & cantidad de sustancia & mol\\
 & & intensidad luminosa & cd\\
\end{tabular}\par\vspace{0.25em}
\lbl{Prefijos}\par
\begin{tabular*}{\linewidth}{@{\extracolsep{\fill}}cccccccc@{}}
G & M & k & d & c & m & \textmu & n\\
$10^{9}$ & $10^{6}$ & $10^{3}$ & $10^{-1}$ & $10^{-2}$ & $10^{-3}$ & $10^{-6}$ & $10^{-9}$\\
\end{tabular*}\par\vspace{0.25em}
\lbl{Unidades derivadas}\par
\begin{tabularx}{\linewidth}{@{}Xll@{}}
\toprule
Magnitud & Unidad & En unidades base\\
\midrule
fuerza & N & $\mathrm{kg\,m/s^2}$\\
trabajo, energ\'ia & $\mathrm{J}=\mathrm{N\cdot m}$ & $\mathrm{kg\,m^2/s^2}$\\
potencia & $\mathrm{W}=\mathrm{J/s}$ & $\mathrm{kg\,m^2/s^3}$\\
presi\'on & $\mathrm{Pa}=\mathrm{N/m^2}$ & $\mathrm{kg/(m\,s^2)}$\\
velocidad & $\mathrm{m/s}$ & \\
aceleraci\'on & $\mathrm{m/s^2}$ & \\
momento de una fuerza (2.\textordmasculine) & $\mathrm{N\cdot m}$ & no se escribe J\\
\bottomrule
\end{tabularx}\par\vspace{0.2em}
\lbl{Sin unidades} lo que va dentro de $\sin$, $\cos$, $e^{x}$: $\omega t$ en radianes,
$e^{-t/\tau}$ con $\tau$ en s.\par
\lbl{Convertir} multiplica por uno:
$72{,}0\units{km/h}\cdot\dfrac{1000\units{m}}{1\units{km}}\cdot\dfrac{1\units{h}}{3600\units{s}}=20{,}0\mps$.
\end{unitbox}

\begin{measbox}{9 \textperiodcentered\ Cifras significativas e incertidumbre}
\begin{tabularx}{\linewidth}{@{}Xrc@{}}
\toprule
Regla & Ejemplo & Cifras\\
\midrule
D\'igitos $\ne0$ y ceros entre ellos: cuentan & $305$ & 3\\
Ceros iniciales: no cuentan & $0{,}00450$ & 3\\
Ceros finales tras la coma: cuentan & $12{,}30$ & 4\\
Ceros finales de un entero: ambiguos & $1200$ & ?\\
\bottomrule
\end{tabularx}\par
$1200$: escribe $1{,}2\cdot10^3$ (2 cifras) o $1{,}200\cdot10^3$ (4 cifras).\par\vspace{0.15em}
\lbl{Sumas y restas} menos decimales: $12{,}3+0{,}45=12{,}75\to12{,}8$\par
\lbl{Productos y cocientes} menos cifras: $2{,}4\cdot3{,}12=7{,}488\to7{,}5$\par
\lbl{Resultado} $\Delta x$ a 1 cifra (2 si empieza por 1); el valor, al mismo decimal:
$9{,}8137\pm0{,}0462\to(9{,}81\pm0{,}05)\mpss$\par
\lbl{Propagar} sumas: absolutas; productos y cocientes: relativas, $\Delta x/|x|$.
\vspace{-0.1em}
\[
\Delta(a\pm b)\approx\Delta a+\Delta b\qquad
\frac{\Delta(ab)}{|ab|}\approx\frac{\Delta(a/b)}{|a/b|}\approx\frac{\Delta a}{|a|}+\frac{\Delta b}{|b|}
\]
\end{measbox}
\end{minipage}
\end{document}
